% A.39 is reserved by the root for the exact cubic-quotient obstruction.
\subsection{The cubic quotient recognizes the first two mixed-label cases}

\begin{statement}{A.39}{Proved}{The realizations $\lambda_0=6$ and $7$ are recognized}
Every noncuspidal Selmer section in A.14 whose labeled integral
Legendre realization in A.16 has parameter $\lambda_0=6$ or $7$ is
a rational point section.  Together with A.36--A.37, this recognizes
every such realization with integral parameter $2\leq\lambda_0\leq9$.
\end{statement}
\begin{proof}
We first construct the exact finite quotient used in the proof.
For a two-dimensional $\F_3$-vector space $M$, let
\[
 \mathcal P(M)=\left\{\text{partitions of the four elements of
 $\mathbf P(M)$ into two unordered pairs}\right\}.
\]
This is a three-element set with its functorial
$\operatorname{GL}_2(\F_3)$-action.  The map on the four projective
lines has kernel $\{\pm I\}$ and identifies
$\operatorname{PGL}_2(\F_3)$ with $S_4$: it is injective, and both
groups have order 24.  Acting on the three pair-partitions gives
the explicit quotient $S_4\to S_3$, whose kernel consists of the
identity and the three double transpositions.  Thus
$\operatorname{GL}_2(\F_3)\to\operatorname{Sym}(\mathcal P(M))$
is onto, with kernel of order eight.  In particular, an isomorphism
of the original Galois modules gives an isomorphism of these
three-element Galois sets, without a choice of symplectic multiplier.

For $E_t:y^2=x(x-1)(x-t)$ the four projective lines of $E_t[3]$
are canonically identified with the four roots $r_1,\ldots,r_4$ of
\[
            3x^4-4(1+t)x^3+6tx^2-t^2.
\]
Indeed, a line has two nonzero points $P,-P$, with the same
$x$-coordinate; these are exactly the roots of the 3-division
polynomial.  Send a partition $\{\{i,j\},\{k,l\}\}$ to
$z=r_ir_j+r_kr_l$.  The elementary symmetric functions of the four
roots are $4(1+t)/3,2t,0,-t^2/3$.  Substitution in the product of
the three factors $z-r_ir_j-r_kr_l$ gives
\[
 R_t(z)=\left(z-\frac{2t}{3}\right)^3+
                    \frac{16t^2(1-t)^2}{27}.
\tag{A.39.1}
\]
This polynomial has distinct roots on $U$, so the displayed map
identifies the three-element sets.  Put $D=4t(1-t)$ and
$w=3z-2t$.  Equation (A.39.1) becomes $w^3=-D^2$.  Since $D$ and
$w$ are units on this cover, the actual inverse coordinate changes
\[
                 v=-D/w,\qquad w=-v^2
\tag{A.39.2}
\]
identify it with the finite etale cover
\[
                         v^3=4t(1-t).
\tag{A.39.3}
\]
The polynomial computation, the order-eight kernel, and the action
of $U(1),L(1),\operatorname{diag}(-1,1)$ on the three partitions
are independently checked in \texttt{code/legendre\_checks.py}.

Let $A,B\in\widehat{\Q^\times}$ be the pullbacks of $t,1-t$ by the
section.  Formula (A.16.3) and its labeled integral comparison give
\[
     A\equiv\lambda_0,\qquad B\equiv1-\lambda_0
                      \pmod{(\Q^\times)^2}.
\tag{A.39.4}
\]
The finite Galois set obtained by pulling back (A.39.3) has Kummer
class $\kappa_3(4)+\kappa_3(A)+\kappa_3(B)$.  The mod-three part
of A.16, followed by the functor $\mathcal P$, identifies this set
with the one for $E_{\lambda_0}$, namely the roots of
\[
                      v^3=4\lambda_0(1-\lambda_0).
\tag{A.39.5}
\]
We will contradict this identification at one specified inertia
group if a section with either realization were not a point.

For $\lambda_0=6$, the potentially multiplicative primes are exactly
$3,5$; for $\lambda_0=7$, they are exactly $3,7$.  This follows
directly from
\[
          j(6)=1906624/225,\qquad j(7)=5088448/441.
\]
Both have $2$-adic $j$-valuation six.  By A.37 no other finite prime
can be a local cusp place.  The same result recognizes every section
having at most one cusp label.  Hence a section not already
recognized must be cuspidal at both listed primes, with distinct
labels.  The local valuations in (A.39.4) force those labels:
\[
\begin{array}{c|cc}
 \lambda_0&\text{cusp at }3&\text{other cusp}\\ \hline
 6&0&1\text{ at }5\\
 7&1&0\text{ at }7.
\end{array}
\tag{A.39.6}
\]
For example, at $3$ in the first row the valuation of $A$ is odd
and that of $B$ is even.  Cusp 1 would make $A=1$ locally, whereas
cusp infinity would make $A/B=-1$ locally; both violate those
valuation parities.  The other three entries follow by the same
valuation check, with $A$ and $B$ interchanged where appropriate.

The section at 2 is an ordinary point $x_2$.  Write
$a=v_2(x_2)$, $b=v_2(1-x_2)$, honest integers.  In the first row of
(A.39.6), (A.39.4) says that $a$ is odd and $b$ is even.  If $a<0$
then $b=a$, a contradiction; hence $a>0$ and $b=0$.  In the second
row $a$ is even and $b$ is odd, so the identical argument with the
two coordinates exchanged gives $a=0$, $b>0$.

At every odd prime outside the two indicated bad primes, the
realized curve has good reduction.  The local integral comparison
and the Neron--Ogg--Shafarevich criterion give good reduction for
the local Legendre curve.  At an odd prime a nonzero valuation of
$x_p$ or $1-x_p$ would give negative $j$-valuation by the displayed
Legendre formula in A.14.  Thus both valuations there are zero.
At a cusp 0 the regular unit $1-t$ has value 1, and at a cusp 1
the regular unit $t$ has value 1.  Combining these facts with the
calculation at 2 and the signs in (A.39.4) gives the exact global
completed classes
\[
 \begin{array}{c|c}
   \lambda_0=6&B=-5^\beta\quad(\beta\in\widehat\Z),\\
   \lambda_0=7&A=7^\alpha\quad(\alpha\in\widehat\Z).
 \end{array}
\tag{A.39.7}
\]
Here there is no unrecorded unit factor: the homomorphism given by
the sign and all prime valuations is injective on
$\widehat{\Q^\times}$, as is seen at each finite power quotient.
With just one nonzero prime-valuation coordinate it gives exactly
the stated power of that prime.

At $3$ the class in the relevant row of (A.39.7) is 1 by (A.39.6).
The cubes of units modulo 9 are precisely $1,8$.  Neither $5$ nor
$7$ is a cube in $\Q_3$, and $-1$ itself is a cube.  Reduction of
the completed class modulo cubes therefore gives
\[
                \beta\equiv0\pmod3\quad\text{or}\quad
                \alpha\equiv0\pmod3,
\tag{A.39.8}
\]
respectively.  The relevant class $B$, or $A$, is consequently a
cube globally in the finite Kummer quotient.

At the other cusp place $q=5$ or $q=7$, the remaining coordinate is
1 by (A.39.6).  Thus the local class of $4AB$ is the unit class of
4 modulo cubes.  Since $q\ne3$, the roots of a unit cube equation
form an unramified Galois set: the algebra
$\Z_q[v]/(v^3-u)$ is finite etale for $u\in\Z_q^\times$.
For the realized curve, however, (A.39.5) has right-hand side
\[
                 -120\quad\text{at }q=5,\qquad
                 -168\quad\text{at }q=7.
\]
Each has $q$-valuation one.  Its cubic polynomial is Eisenstein,
so the associated degree-three field is totally ramified, and
inertia acts nontrivially on its roots.  This contradicts the
Galois-set identification obtained from the mod-three integral
comparison.  Both realizations are therefore recognized.

Finally, $\lambda_0=2$ is covered by A.36.  The parameters
$3,4,5,8,9$ have respectively the single potentially multiplicative
prime $3,3,5,7,3$, as follows by factoring the denominator of their
$j$-invariant.  A.37 recognizes them.  This proves the last
assertion without claiming a classification for the remaining
rational parameters.
\end{proof}
